\section{Two overlapping sweeps in a tetrahedron}
\label{sec:tetrahedron}

We now choose two sweeps for which the weighted cost is strictly
smaller than the half-area bound. Write physical points as
$(x,y,ht)$, where $h=\sqrt2$, and let
\begin{equation}\label{eq:tetrahedron}
 K=\{(x,y,ht):|x|\le1+t,\ |y|\le1-t\}.
\end{equation}
The inequalities imply $-1\le t\le1$. The four vertices are
$(\pm2,0,h)$ and $(0,\pm2,-h)$, and their pairwise distances
are four. This is a regular tetrahedron.

\subsection{The reference area}

For a polytope with face areas $S_j$, outward unit normals $n_j$,
and a unit vector $u$, Cauchy's projection formula takes the form
\begin{equation}\label{eq:cauchy}
 \area(\pi_{u^\perp}K)=\frac12\sum_j S_j|n_j\cdot u|.
\end{equation}
See Martini \cite[p.~83, equation~(1)]{Martini1991} for the
projection-body formulation. The elementary polytope proof is to
count the two endpoints of almost every line parallel to $u$:
projecting all faces counts its shadow twice, whereas edges and
faces parallel to $u$ contribute zero area.

Each face of \eqref{eq:tetrahedron} has area $4\sqrt3$, and its
normal is one of the normalizations of
$(\pm1,0,-1/h)$ and $(0,\pm1,1/h)$.
For a unit vector $u=(u_1,u_2,u_3)$, formula~\eqref{eq:cauchy}
and $|a+b|+|a-b|=2\max(|a|,|b|)$ give
\[
 \area(\pi_{u^\perp}K)
 =4h\left(\max\{|u_1|,|u_3|/h\}
          +\max\{|u_2|,|u_3|/h\}\right).
\]
If the two maxima are $X,Y$, then
$X^2+Y^2\ge u_1^2+u_2^2$ and $XY\ge u_3^2/h^2$.
Since $h^2=2$, it follows that $(X+Y)^2\ge1$.
Equality is attained at $u=(1,0,0)$, so
\begin{equation}\label{eq:min-area}
 A_{\min}(K)=4h.
\end{equation}
The upper half of $K$ is covered by an $x$-parallel cylinder of
base area $h$, and the lower half by a $y$-parallel cylinder of
the same area. Their total cost is $2h=A_{\min}(K)/2$.

\subsection{The perturbed segments cover the whole tetrahedron}

For $q\in[-1,1]$, define
\[
 G(q)=\frac{1+q^2}{2},\qquad H(q)=\frac{q^2+q^4}{2},\qquad
 R_\eps(q)=1+\eps^2\bigl(H(q)+5\eps\bigr),\quad \eps>0.
\]
The upper sweep is
\begin{equation}\label{eq:upper-sweep}
 U_\eps=\{(x,\rho q-\eps xG(q),
                     h(1-\rho+\eps xq)):
       |x|\le2,\ |q|\le1,\ 0\le\rho\le R_\eps(q)\}.
\end{equation}
Its intercept ray starts at $(0,0,h)$ in direction $(0,q,-h)$.
The identity $G'(q)=q$ makes the derivative of its segment direction
$(1,-\eps G(q),h\eps q)$ equal to $-\eps(0,q,-h)$, exactly
the radial alignment required by Theorem~\ref{thm:radial}.
The constant term in $G$ gives
$G(q)-q^2=(1-q^2)/2$, which vanishes at both angular endpoints.
This will keep the endpoint levels unchanged when we select a line
through a point of $K$.

The sweep is compact. Let $\mathcal R(x,y,z)=(y,-x,-z)$, an orthogonal
map preserving $K$. Its image $\mathcal R U_\eps$ is the lower
sweep, which can equivalently be written as
\begin{equation}\label{eq:lower-sweep}
 \{(\sigma p+\eps yG(p),y,h(-1+\sigma+\eps yp)):
       |y|\le2,\ |p|\le1,\ 0\le\sigma\le R_\eps(p)\}.
\end{equation}
Indeed, apply $\mathcal R$ in \eqref{eq:upper-sweep} and replace
its free coordinate $x$ by $-y$. We will show that the two sweeps
overlap enough to cover $K$.

\begin{lemma}\label{lem:crossing}
If $F\in C^1([a,b])$, $a<b$, and $F(a)\le z\le F(b)$,
there is $r\in[a,b]$ with $F(r)=z$ and $F'(r)\ge0$,
using one-sided derivatives at the endpoints.
\end{lemma}
\begin{proof}
If $F(a)=z$ and $F'(a)\ge0$, use $r=a$.
Otherwise $F$ is below $z$ either at $a$ or immediately to its
right. Starting at such a point, take the first subsequent point
where $F=z$; it exists because $F(b)\ge z$. Its left difference
quotients are nonnegative, so its derivative is nonnegative.
\end{proof}

\begin{proposition}\label{prop:coverage}
For every $\eps>0$, the compact sets $U_\eps$ and
$\mathcal R U_\eps$ cover $K$.
\end{proposition}
\begin{proof}
Fix $(x,y,ht)\in K$. Applying Lemma~\ref{lem:crossing} twice
gives $q,p\in[-1,1]$ satisfying
\begin{align}
 y&=(1-t)q-\eps x(1-q^2)/2,
 &\rho&=1-t+\eps xq\ge0,\label{eq:root-one}\\
 x&=(1+t)p+\eps y(1-p^2)/2,
 &\sigma&=1+t-\eps yp\ge0.\label{eq:root-two}
\end{align}
In each equation the endpoint values bracket the coordinate on
the left, and the derivative is precisely the indicated radial
parameter. This includes zero-height slices and boundary points.
Rearranging gives the segment representations
\eqref{eq:upper-sweep} and \eqref{eq:lower-sweep}; only their
upper radial bounds remain to be checked.

If both bounds failed, then
\begin{equation}\label{eq:joint-failure}
 \eps yp+\eps^2(H(p)+5\eps)<t
       <\eps xq-\eps^2(H(q)+5\eps).
\end{equation}
As $|x|,|y|\le2$ and both cap corrections are nonnegative,
$|t|\le2\eps$. Equations~\eqref{eq:root-one}--\eqref{eq:root-two}
then imply
\[
 |x-p|\le3\eps,\qquad |y-q|\le3\eps.
\]
Moreover, \eqref{eq:joint-failure} puts $t/\eps$ between
$yp$ and $xq$. Both differ from $pq$ by at most $3\eps$, so
\begin{equation}\label{eq:interface-error}
 |t/\eps-pq|\le3\eps.
\end{equation}

Put $D(r)=(1-r^2)/2$, so $0\le D(r)\le1/2$ on $[-1,1]$.
The root equations give the exact identity
\[
 \frac{xq-yp}{\eps}
   =2pq\frac t\eps+yqD(p)+xpD(q).
\]
Replacing $t/\eps$ by $pq$ costs at most $6\eps$, by
\eqref{eq:interface-error}. Replacing $y$ by $q$ and $x$ by
$p$ costs at most $3\eps$ in total. Hence
\begin{align*}
 \frac{xq-yp}{\eps}
 &\le2p^2q^2+q^2D(p)+p^2D(q)+9\eps\\
 &=\frac{p^2+q^2}{2}+p^2q^2+9\eps.
\end{align*}
This is why the cap contains $H(q)=q^2G(q)$: the inequality
$2p^2q^2\le p^4+q^4$ gives
\[
 \frac{p^2+q^2}{2}+p^2q^2\le H(p)+H(q).
\]
The two cubic cap corrections contribute $10\eps$ after division
by $\eps^2$, exceeding the error $9\eps$.
Indeed, \eqref{eq:joint-failure} requires $(xq-yp)/\eps$ to be
strictly greater than $H(p)+H(q)+10\eps$. This contradiction
proves the proposition. All estimates were uniform over the closed
tetrahedron.
\end{proof}

\subsection{Finite bases and the weighted cost}

We now hold $\eps$ fixed. In the intercept plane $P=\{x=0\}$,
the upper sweep has the data
\[
 p_0=(0,0,h),\quad e=(1,0,0),\quad
 a=(0,0,-h),\quad b=(0,1,0),\quad L=2,
\]
and
\[
 w(q)=(0,-\eps G(q),h\eps q),\qquad
 w'(q)=-\eps(a+qb).
\]
Thus the hypothesis of Theorem~\ref{thm:radial} holds with
$\lambda=-\eps$, $\kappa=h$, and radial cap $R_\eps$.
Also
\[
 \norm{e+w(q)}^2=1+\eps^2D_0(q),\qquad
 D_0(q)=G(q)^2+2q^2.
\]
For every partition into $N$ intervals, the theorem and the orthogonal
map $\mathcal R$ cover $K$ by $2N$ cylinders with nondegenerate compact
triangular bases. As the mesh tends to zero, their total areas converge to
\begin{equation}\label{eq:cost-integral}
 2hI(\eps),\qquad
 I(\eps)=\int_{-1}^1
       \frac{R_\eps(q)^2}{2\sqrt{1+\eps^2D_0(q)}}\,dq.
\end{equation}
Taking actual swept projections instead gives compact bases with
the same limiting total area, for any choice of freezing tags.
Their finite areas may differ from those of the triangles.

\subsection{The saving and the final finite choice}

It remains to put the limiting cost strictly below $2h$.
The expansion is uniform for $q\in[-1,1]$:
\[
 \frac{R_\eps(q)^2}{2\sqrt{1+\eps^2D_0(q)}}
 =\frac12+\eps^2\left(H(q)-\frac{D_0(q)}4\right)
       +5\eps^3+O(\eps^4).
\]
Direct integration gives
\[
 \int_{-1}^1H(q)\,dq=\frac8{15},\qquad
 \int_{-1}^1\frac{D_0(q)}4\,dq=\frac{17}{30}.
\]
Therefore
\begin{equation}\label{eq:strict-expansion}
 I(\eps)=1-\frac{\eps^2}{30}+10\eps^3+O(\eps^4)<1
\end{equation}
for all sufficiently small positive $\eps$.
The cubic cap padding that guaranteed overlap is smaller than
the quadratic projection saving.

Choose one such $\eps$ and put
$\Delta=2h(1-I(\eps))>0$.
For the triangular covers, the limit \eqref{eq:cost-integral}
gives a sufficiently fine finite partition, with $N$ intervals, whose
cost is less than $2hI(\eps)+\Delta/2<2h$.
Proposition~\ref{prop:coverage} and Theorem~\ref{thm:radial}
give coverage of every point of $K$ for this same $N$.
The actual swept projections also give finite covers within this gap
by their convergence to the same limit.

By \eqref{eq:min-area}, the triangular cover has cost strictly
below $A_{\min}(K)/2$. A Euclidean similarity transports it to
every regular tetrahedron of positive edge length, preserving
triangularity and multiplying both areas by the same factor.
Finally each individual projection area is at least $A_{\min}(K)$.
For the axes and bases of any one of these finite covers,
\[
 \sum_i\frac{\area(B_i)}{\area(\pi_{u_i^\perp}K)}
 \le\frac{\sum_i\area(B_i)}{A_{\min}(K)}<\frac12.
\]
This proves the two negative resolutions stated in the introduction.

Explicit polygonal alternatives with the same limiting cost are given
in Appendix~\ref{sec:finite-footprints}.
